% DSCT Week 4 Decision Gate — reasoning write-up template
% Compile with the Week 3 pinned container:
%   podman run --rm -v "$PWD":/work -w /work \
%     ghcr.io/jeremy-evert/dsct-week3-latex:latest \
%     latexmk -pdf decision-gate-template.tex
% (or any local LaTeX: pdflatex / latexmk / tectonic)
\documentclass[11pt]{article}
\usepackage[margin=1in]{geometry}
\usepackage{amsmath, amssymb, amsthm}
\usepackage{enumitem}
\usepackage{parskip}
\usepackage[hidelinks]{hyperref}

\newtheorem*{claim}{Claim}
\newtheorem*{counterexample}{Counterexample}

\title{Week 4 Decision Gate: \emph{King Octopus and the servants}}
\author{Jeremy Evert}
\date{\today}

\begin{document}
\maketitle

\section*{Source and Problem Statement} The following puzzle is from Puzzle Prime: \begin{center} \href{https://www.puzzleprime.com/puzzles/brain-teasers/deduction/king-octopus-and-his-servants/} {\textbf{King Octopus and His Servants}} \end{center} King Octopus has servants with 6, 7, or 8 legs. The servants with 7 legs always lie, and the servants with 6 or 8 legs always tell the truth. One day four of the king's servants had the following conversation: \begin{quote} ``We together have 28 legs,'' said the first octopus. ``We together have 27 legs,'' said the second octopus. ``We together have 26 legs,'' said the third octopus. ``We together have 25 legs,'' said the fourth octopus. \end{quote} \textbf{Question:} Which of the four servants told the truth? \medskip \noindent \textit{Source:} \href{https://www.puzzleprime.com/puzzles/brain-teasers/deduction/king-octopus-and-his-servants/} {Puzzle Prime, ``King Octopus and His Servants.''}

\section*{Rules / Assumptions}
What the key terms mean here: \emph{valid}, \emph{follows}, ``for every
integer'', a servant has exactly one of $6,7,8$ legs, \dots

\section*{Work}
% Translate the claim into precise form, then reason.
\begin{claim}
State the claim precisely.
\end{claim}

\begin{proof}
Steps, each justified by a definition or a prior step. Or, if the claim is
false:
\end{proof}

\begin{counterexample}
An object in the claim's domain that satisfies the hypothesis and violates
the conclusion. Show the values.
\end{counterexample}

\section*{Check Your Answer}
Either verify \emph{every} proof step against a definition, or verify the
counterexample really satisfies the hypothesis and violates the conclusion.
For an argument-validity problem, exhibit premises that are all true while
the conclusion is false (or argue that is impossible).

\section*{One-Sentence Summary}
Exactly what you established --- and, where it matters, distinguish
``the conclusion happens to be true'' from ``the argument is valid.''

\section*{Strongest disagreement}
The best objection a sharp, skeptical reader would raise, stated honestly,
and whether it survives.

\end{document}
